% GENERATED FILE. Edit canonical lessons, then run npm run generate:books.
% canonical-source-hash: fnv1a64:5d483e27590ceeb9
% canonical-lessons: TA-C259-patam-etu, TA-C259-nati, TA-C259-raci, TA-C259-palutu-par, TA-C259-olunkupatuttu

\chapter{To take a photo, To act, To enjoy, To repair, To tidy up}
\label{ch:ta-to-take-a-photo-to-act-to-enjoy-to-repair-to-tidy-up}
\hlchaptermodality{259}

\begin{chapteropening}
\textbf{By the end of this chapter:} \emph{I can say to take a photo, to act, to enjoy, to repair and to tidy up.}

Complete the last lesson of chapter 259: I can say to take a photo, to act, to enjoy, to repair and to tidy up.
\end{chapteropening}

\section[\texorpdfstring{\ta{படம்} \ta{எடு} (paṭam eṭu)}{படம் எடு (paṭam eṭu)}]{\ta{படம்} \ta{எடு} (paṭam eṭu) — to take a photo}

\label{lesson:TA-C259-patam-etu}



\begin{quote}
Before the new one: say the Tamil for to comfort, then the Tamil for to shake hands.
\end{quote}

\subsection*{You'll want to know: \ta{படம்} \ta{எடு}}
\textbf{\ta{படம்} \ta{எடு}} — \emph{paṭam eṭu} — \textquotedblleft{}to take a photo\textquotedblright{}.

\begin{grammarlens}[title={What you've built}]
One more word for this chapter.
\end{grammarlens}

\subsection*{Your turn}
\begin{itemize}
\raggedright
  \item \emph{Say it:} \emph{paṭam eṭu}
  \item \emph{Say it:} \emph{paṭam eṭu}, once more
  \item \emph{Say it:} say \emph{paṭam eṭu}
  \item \emph{Recall:} say the Tamil for to comfort, then the Tamil for to shake hands, then say \emph{paṭam eṭu} again
\end{itemize}

\begin{tcolorbox}[breakable,colback=teal!4,colframe=teal!35!black,title={Before you move on}]
What does \ta{படம்} \ta{எடு} mean? (\textquotedblleft{}To take a photo\textquotedblright{}.) Say it once more.
\end{tcolorbox}
\section[\texorpdfstring{\ta{நடி} (naṭi)}{நடி (naṭi)}]{\ta{நடி} (naṭi) — to act (in a play or film)}

\label{lesson:TA-C259-nati}



\begin{quote}
Before the new one: say the Tamil for to shake hands, then the Tamil for to take a photo.
\end{quote}

\subsection*{You'll want to know: \ta{நடி}}
\textbf{\ta{நடி}} — \emph{naṭi} — \textquotedblleft{}to act (in a play or film)\textquotedblright{}.

\begin{grammarlens}[title={What you've built}]
One more word for this chapter.
\end{grammarlens}

\subsection*{Your turn}
\begin{itemize}
\raggedright
  \item \emph{Say it:} \emph{naṭi}
  \item \emph{Say it:} \emph{naṭi}, once more
  \item \emph{Say it:} say \emph{naṭi}
  \item \emph{Recall:} say the Tamil for to shake hands, then the Tamil for to take a photo, then say \emph{naṭi} again
\end{itemize}

\begin{tcolorbox}[breakable,colback=teal!4,colframe=teal!35!black,title={Before you move on}]
What does \ta{நடி} mean? (\textquotedblleft{}To act (in a play or film)\textquotedblright{}.) Say it once more.
\end{tcolorbox}
\section[\texorpdfstring{\ta{ரசி} (raci)}{ரசி (raci)}]{\ta{ரசி} (raci) — to enjoy, to appreciate}

\label{lesson:TA-C259-raci}



\begin{quote}
Before the new one: say the Tamil for to take a photo, then the Tamil for to act.
\end{quote}

\subsection*{You'll want to know: \ta{ரசி}}
\textbf{\ta{ரசி}} — \emph{raci} — \textquotedblleft{}to enjoy, to appreciate\textquotedblright{}.

\begin{grammarlens}[title={What you've built}]
One more word for this chapter.
\end{grammarlens}

\subsection*{Your turn}
\begin{itemize}
\raggedright
  \item \emph{Say it:} \emph{raci}
  \item \emph{Say it:} \emph{raci}, once more
  \item \emph{Say it:} say \emph{raci}
  \item \emph{Recall:} say the Tamil for to take a photo, then the Tamil for to act, then say \emph{raci} again
\end{itemize}

\begin{tcolorbox}[breakable,colback=teal!4,colframe=teal!35!black,title={Before you move on}]
What does \ta{ரசி} mean? (\textquotedblleft{}To enjoy, to appreciate\textquotedblright{}.) Say it once more.
\end{tcolorbox}
\section[\texorpdfstring{\ta{பழுது} \ta{பார்} (paḻutu pār)}{பழுது பார் (paḻutu pār)}]{\ta{பழுது} \ta{பார்} (paḻutu pār) — to repair}

\label{lesson:TA-C259-palutu-par}



\begin{quote}
Before the new one: say the Tamil for to act, then the Tamil for to enjoy.
\end{quote}

\subsection*{You'll want to know: \ta{பழுது} \ta{பார்}}
\textbf{\ta{பழுது} \ta{பார்}} — \emph{paḻutu pār} — \textquotedblleft{}to repair\textquotedblright{}.

\begin{grammarlens}[title={What you've built}]
One more word for this chapter.
\end{grammarlens}

\subsection*{Your turn}
\begin{itemize}
\raggedright
  \item \emph{Say it:} \emph{paḻutu pār}
  \item \emph{Say it:} \emph{paḻutu pār}, once more
  \item \emph{Say it:} say \emph{paḻutu pār}
  \item \emph{Recall:} say the Tamil for to act, then the Tamil for to enjoy, then say \emph{paḻutu pār} again
\end{itemize}

\begin{tcolorbox}[breakable,colback=teal!4,colframe=teal!35!black,title={Before you move on}]
What does \ta{பழுது} \ta{பார்} mean? (\textquotedblleft{}To repair\textquotedblright{}.) Say it once more.
\end{tcolorbox}
\section[\texorpdfstring{\ta{ஒழுங்குபடுத்து} (oḻuṅkupaṭuttu)}{ஒழுங்குபடுத்து (oḻuṅkupaṭuttu)}]{\ta{ஒழுங்குபடுத்து} (oḻuṅkupaṭuttu) — to tidy up, to arrange}

\label{lesson:TA-C259-olunkupatuttu}



\begin{quote}
Before the new one: say the Tamil for to enjoy, then the Tamil for to repair.
\end{quote}

\subsection*{You'll want to know: \ta{ஒழுங்குபடுத்து}}
\textbf{\ta{ஒழுங்குபடுத்து}} — \emph{oḻuṅkupaṭuttu} — \textquotedblleft{}to tidy up, to arrange\textquotedblright{}.

\begin{grammarlens}[title={What you've built}]
One more word for this chapter.
\end{grammarlens}

\subsection*{Your turn}
\begin{itemize}
\raggedright
  \item \emph{Say it:} \emph{oḻuṅkupaṭuttu}
  \item \emph{Say it:} \emph{oḻuṅkupaṭuttu}, once more
  \item \emph{Say it:} say \emph{oḻuṅkupaṭuttu}
  \item \emph{Recall:} say the Tamil for to enjoy, then the Tamil for to repair, then say \emph{oḻuṅkupaṭuttu} again
\end{itemize}

\begin{tcolorbox}[breakable,colback=teal!4,colframe=teal!35!black,title={Before you move on}]
What does \ta{ஒழுங்குபடுத்து} mean? (\textquotedblleft{}To tidy up, to arrange\textquotedblright{}.) Say it once more.
\end{tcolorbox}
